\section{Hướng dẫn cấu hình Router và các thiết bị}
\label{sec:configuration-guide}

Các bước dưới đây là quy trình khái quát; tên menu và khả năng tùy chọn phụ thuộc
nhà sản xuất/model. Trước khi thao tác, ghi nhận model, phiên bản firmware, sơ đồ
kết nối và trạng thái ban đầu. Không đăng tải thông tin quản trị hoặc mật khẩu.

\subsection{Cấu hình Router/điểm truy cập}
\begin{enumerate}
  \item Kết nối máy quản trị với cổng LAN hoặc giao diện quản lý an toàn; đăng
    nhập bằng thông tin hợp lệ và đổi mật khẩu quản trị mặc định.
  \item Cập nhật firmware từ nguồn chính thức; sao lưu cấu hình hiện hành nếu
    thiết bị hỗ trợ.
  \item Cấu hình kết nối WAN theo thông số nhà cung cấp dịch vụ; đặt địa chỉ
    LAN, mặt nạ mạng và DHCP phù hợp, tránh trùng dải địa chỉ đang dùng.
  \item Bật Wi-Fi cho băng tần cần sử dụng; đặt SSID, chọn kênh/kênh rộng phù
    hợp môi trường và cấu hình WPA2-AES hoặc WPA3-Personal với mật khẩu mạnh.
  \item Tắt WPS và dịch vụ quản trị từ WAN nếu không cần; lưu cấu hình, khởi
    động lại khi thiết bị yêu cầu và ghi lại các giá trị đã thiết lập.
\end{enumerate}

\subsection{Kết nối Laptop, Desktop và Smartphone}
Trên mỗi thiết bị, bật Wi-Fi, chọn đúng SSID và xác thực; với Desktop dùng cáp,
kết nối vào cổng LAN. Kiểm tra địa chỉ IP, gateway và DNS nhận được, sau đó
kiểm tra lần lượt khả năng liên lạc với gateway, thiết bị cùng mạng và Internet.
Ghi kết quả thực tế vào phần thử nghiệm; không điền kết quả dự kiến như kết quả
đo được.

\subsection{Thiết lập và kiểm chứng mạng Ad-hoc (IBSS) trên nền Linux}
\label{sec:adhoc-config}

Trước khi cấu hình, cần xác nhận nền tảng có hỗ trợ chế độ ngang hàng hay không.
Trên Windows~11, lệnh \texttt{netsh wlan show drivers} trả về dòng
\emph{Hosted network supported: No}, nghĩa là hệ điều hành không còn hỗ trợ chế
độ Ad-hoc (minh chứng tại Hình~\ref{fig:win-hosted}); Android và macOS cũng
không hỗ trợ. Vì vậy quy trình dưới đây được thực hiện trên Linux, sử dụng module
\texttt{mac80211\_hwsim} để mô phỏng phần cứng Wi-Fi và hai network namespace để
tách hai nút thành hai máy độc lập. Không được gọi nhầm chế độ hotspot hay
Wi-Fi Direct là \ac{ibss}.

\begin{enumerate}[leftmargin=*]
  \item \textbf{Kiểm tra nền tảng và tạo card Wi-Fi ảo.} Nạp module mô phỏng với
    hai radio: \texttt{sudo modprobe mac80211\_hwsim radios=2}, sau đó kiểm tra
    bằng \texttt{iw dev}. Mỗi radio tương ứng một giao diện (\texttt{wlan0},
    \texttt{wlan1}) có kênh, công suất phát và thuộc tính phần cứng riêng.
    Giao diện \texttt{P2P-device} do trình điều khiển tự tạo không ảnh hưởng
    đến kết nối Ad-hoc.

  \item \textbf{Tách mỗi radio vào một network namespace.} Nếu hai giao diện nằm
    chung một hệ thống, chúng dùng chung bảng định tuyến nên Linux coi cả hai
    địa chỉ IP là của chính nó và chuyển gói qua loopback, khiến lưu lượng không
    đi qua ``sóng'' vô tuyến ảo. Thực hiện:
    \texttt{sudo ip netns add ns1} và \texttt{sudo ip netns add ns2};
    \texttt{sudo iw phy phy0 set netns name ns1};
    \texttt{sudo iw phy phy1 set netns name ns2}. Kiểm tra lại bằng
    \texttt{sudo ip netns exec ns1 iw dev} và
    \texttt{sudo ip netns exec ns2 iw dev}.

  \item \textbf{Chuyển giao diện sang chế độ \ac{ibss} và tham gia mạng.} Với
    mỗi nút, hạ giao diện xuống trước khi đổi chế độ (nếu không, nhân hệ điều
    hành báo lỗi \emph{Device or resource busy}). Với Node~1 trong \texttt{ns1}:
    \texttt{sudo ip netns exec ns1 ip link set wlan0 down};
    \texttt{sudo ip netns exec ns1 iw dev wlan0 set type ibss};
    \texttt{sudo ip netns exec ns1 ip link set wlan0 up};
    \texttt{sudo ip netns exec ns1 iw dev wlan0 ibss join demoAdhoc 2437};
    \texttt{sudo ip netns exec ns1 ip addr add 10.0.0.1/24 dev wlan0}. Tham số
    \texttt{demoAdhoc} là \ac{ssid} của mạng và mọi nút phải trùng; \texttt{2437}
    là tần số tính bằng MHz, tương ứng kênh~6 của băng tần 2,4~GHz và cũng phải
    trùng giữa các nút. Node~2 thực hiện tương tự với \texttt{ns2},
    \texttt{wlan1} và địa chỉ \texttt{10.0.0.2/24}. Mọi lệnh đều chạy trong
    namespace tương ứng thông qua \texttt{sudo ip netns exec <namespace>}.

  \item \textbf{Kiểm tra liên kết ở tầng 802.11.} Dùng
    \texttt{iw dev wlan0 info} để xác nhận \texttt{type IBSS}, \ac{ssid},
    kênh và tần số; dùng \texttt{iw dev wlan0 station dump} để xem danh sách nút
    lân cận đã phát hiện. Nếu thấy địa chỉ MAC của nút còn lại ở trạng thái
    \emph{authorized} và \emph{associated}, hai nút đã nhận ra nhau.

  \item \textbf{Kiểm tra liên lạc ở tầng mạng.} Chạy
    \texttt{ping -c 4 10.0.0.2} từ Node~1. Gói tin được đóng khung 802.11, truyền
    qua ``sóng ảo'' của \texttt{hwsim} tới giao diện của Node~2. Kết quả đạt yêu
    cầu khi nhận đủ 4/4 phản hồi và tỉ lệ mất gói bằng 0\,\%.

  \item \textbf{Ghi nhận và dọn dẹp.} Lưu ảnh chụp màn hình từng bước vào phần
    minh chứng (Mục~\ref{sec:evidence}) và ghi lại tham số thực tế. Sau khi kết
    thúc, gỡ namespace và module mô phỏng để trả máy về trạng thái ban đầu.
\end{enumerate}

\subsection{Bảng ghi cấu hình thực tế}
\begin{center}
\begin{tabularx}{\linewidth}{X X X}
  \toprule
  Hạng mục & Giá trị thực tế & Minh chứng/ghi chú \\
  \midrule
  Model Router/AP & & \\
  SSID và bảo mật & & \\
  LAN/DHCP & & \\
  Laptop/Desktop/Smartphone & & \\
  Ad-hoc hoặc phương án thay thế & IBSS trên Ubuntu 22.04 LTS (máy ảo
    VMware), \texttt{mac80211\_hwsim radios=2}; \ac{ssid} \texttt{demoAdhoc},
    kênh~6 (2437~MHz), IP \texttt{10.0.0.1/24} và \texttt{10.0.0.2/24}, mạng mở
    & Hình~\ref{fig:hwsim-radios}--\ref{fig:adhoc-ping}; Windows~11 không hỗ trợ
    (Hình~\ref{fig:win-hosted}) \\
  \bottomrule
\end{tabularx}
\end{center}
